\section{Introduction}
\label{sec:intro}

Competition authorities need scalable tools to flag bid-rigging
cartels in public procurement. Existing screens
\citep{imhof2019detecting, huber2019machine} require bid-level
microdata---prices, dispersion measures, distributional
features---that are incomplete or absent in many procurement systems.
This paper proposes a screen that requires only participation records
and win/loss outcomes: \emph{frequent losers} (FL)---firms that never
win yet participate in abnormally many tenders. The screen is
deployable at scale, produces firm-level flags that triage
investigations, and captures information orthogonal to bid-level tools.

The behavioral anomaly is simple: sustained participation with a zero
win rate is difficult to reconcile with profit-maximizing entry when
bidding is costly. It matches patterns documented in prosecuted
cartels, where designated firms submit deliberately losing bids to
simulate competition---what we call \emph{synthetic competition}
\citep{marshall2012economics, asker2010leniency}. A firm that bids on
50 tenders over a decade without winning a single one is either
extraordinarily unlucky or playing a different game.

We validate the screen against external enforcement data. Among
2,735 FL firms identified in S\~ao Paulo's Bolsa Eletr\^onica de
Compras (BEC, 2009--2019), 7.1\% co-participate with
CADE-convicted cartelists---3.5$\times$ the baseline rate
($p < 0.001$). The ROC curve achieves AUC~$= 0.94$, with the
data-driven optimal threshold (Youden's $J = 0.84$) nearly
identical to the baseline IQR rule. A horse-race regression
shows the FL screen captures price information orthogonal to
\citet{imhof2019detecting}-style bid-level features (correlation
0.06), confirming complementarity rather than redundancy.

FL-flagged procurement environments show 3.6--7.7\% higher prices in
within-item, within-year, within-purchasing-unit comparisons. The range
spans temporal cross-fitting (3.6\%), OLS with high-dimensional fixed
effects (6.4\%), and matching (CEM: 7.7\%; IPW: 5.5\%). After
absorbing item fixed effects the raw gap collapses from 2.43 to 0.050
log points (overlap 0.978; Table~\ref{tab:conditional_descstats},
Appendix~\ref{app:competition}), so the regressions compare like with
like. Comparability, however, is not identification.

These are conditional associations, not causal effects. The gap may
reflect collusive environments or unobserved market characteristics
that jointly attract FL firms and higher prices; we cannot fully
separate the two. Sensitivity analysis
(\citealt{cinelli2020making}, $RV = 17.5\%$) bounds the scope for
omitted confounders, and the event-study pre-trends are
presented as an unresolved ambiguity
(Section~\ref{sec:additional_id}). The contribution is primarily
\emph{diagnostic}: FL presence flags procurement environments
systematically associated with higher prices and with known cartel
activity.

We complement the price regressions with a battery of diagnostics
that check whether the screen's behavior is coherent with synthetic
competition. The price effect concentrates in competitive markets;
bid-coordination tests reject exchangeability of FL and non-FL
residuals; FL--winner pairs recur more often than chance predicts;
and the FL premium is larger in convite tenders where the
minimum-bidder rule does not bind (7.6\%) than where it forces
participation. These patterns are individually coherent with cover
bidding but admit non-collusive readings; they serve as supporting
diagnostics, not identifying evidence
(Sections~\ref{sec:results}--\ref{sec:mechanisms}).

\subsection{Contributions}

\emph{First}, we build a firm-level screening tool that uses only
participation frequency and win rates---no bid microdata, no
structural assumptions, no enforcement records. It achieves
AUC~$= 0.94$ against competition-authority convictions and is
deployable at scale.

\emph{Second}, we show that FL-flagged environments have
3.6--7.7\% higher prices in within-item comparisons (a conditional
association, not a causal estimate) and that the screen captures
information orthogonal to existing bid-level tools.

\emph{Third}, we check whether the screen's behavior is coherent with
synthetic competition through independent diagnostics (Bajari--Ye,
dyadic linkage, invitation heterogeneity, CADE co-bidding). The
diagnostics produce coherent results; they do not identify the
mechanism.

We do not claim to identify cover bidders, estimate causal overcharges,
or measure welfare losses.

After describing the data and FL classification
(Sections~\ref{sec:data}--\ref{sec:structural}), we present the
empirical strategy (Section~\ref{sec:empirical}), the CADE
consistency check (Section~\ref{sec:cade}), the price regressions
and mechanism diagnostics
(Sections~\ref{sec:results}--\ref{sec:mechanisms}), robustness
(Section~\ref{sec:robustness}), an explicit discussion of
limitations (Section~\ref{sec:limitations}), and a conclusion
(Section~\ref{sec:conclusion}).
